\section{Material e métodos}
\label{section:method}

O pipeline foi executado em um ambiente contêiner baseado em Ubuntu 24.04. A automação e o gerenciamento das ferramentas foram estruturados via Miniforge3, criando-se o ambiente virtual \cmd{qiime2\_amplicon}. Este ambiente integrou nativamente o QIIME 2 (v. 2024.10) e as dependências em Python 3.10 utilizadas para as análises estatísticas e gráficos (JupyterLab, pandas, seaborn e matplotlib), além dos utilitários de terminal cutadapt, pigz e sra-tools.

Os dados de sequenciamento paired-end foram obtidos a partir do repositório público DNA Databank of Japan, sob o esforço amostral de n = 8 corridas de sequenciamento (IDs: \cmd{DRR030421}-\cmd{DRR030428}). O desenho experimental consistiu em quatro réplicas biológicas para o ambiente artificial controlado (Aquário Churaumi) e quatro réplicas para o ambiente marinho natural (Recife de Corais de Bise, Okinawa, Japão). O download automatizado dos arquivos originais em formato \cmd{.sra} foi realizado via protocolo HTTP, seguido da extração das leituras brutas demultiplexadas em arquivos FASTQ pareados utilizando a ferramenta \cmd{fasterq-dump} (sra-tools). Os arquivos foram compactados com o algoritmo \cmd{pigz} para otimização de armazenamento.

As análises de bioinformática foram centralizadas no ambiente QIIME 2. Inicialmente, um arquivo de manifesto contendo os caminhos absolutos das leituras foi estruturado para importação das sequências brutas pelo formato \cmd{PairedEndFastqManifestPhred33V2}, gerando o artefato base de dados pareados com qualidade Phred-33. A remoção estrita dos primers foi executada com o comando \cmd{qiime cutadapt trim-paired}, que buscou as sequências do marcador universal MiFish na região do gene mitocondrial 12S rRNA (Forward MiFish-U-F: \cmd{5'-GTCGGTAAAACTCGTGCCAGC-3'}) e Reverse (MiFish-U-R: \cmd{5'-CATAGTGGGGTATCTAATCCCAGTTTG-3'}), com correspondência de caracteres curinga. O perfil de qualidade das leituras antes e após a remoção dos primers foi inspecionado visualmente por meio do comando \cmd{qiime demux summarize}.

As leituras limpas foram submetidas ao algoritmo DADA2 via comando \cmd{qiime dada2 denoise-paired} para correção de erros de sequenciamento, filtragem por qualidade, fusão de fitas (paired-end merging) e remoção de quimeras pelo método consensus. Após inspeção da qualidade das corridas, os parâmetros de truncamento foram fixados em 120 pb para as leituras Forward e 115 pb para as Reverse, sem cortes à esquerda (\cmd{trim-left = 0}), garantindo uma zona de sobreposição (overlap) segura de aproximadamente 65 pb para o acoplamento. As Variantes de Sequência de Amplicons (ASVs) resultantes foram tabuladas e classificadas taxonomicamente utilizando um classificador probabilístico Naive Bayes treinado pelo comando \cmd{qiime feature-classifier fit-classifier-naive-bayes}. O treinamento, direcionado para a região do gene 12S amplificada pelos iniciadores, teve como base a biblioteca de referência curada e especializada Mitohelper \cite{Lim.Thompson2021}.

A matriz de abundâncias foi padronizada por rarefação em uma profundidade fixa de 45.000 reads por amostra, valor determinado pela estabilização das curvas e formação do platô inspecionados pelo comando \cmd{qiime diversity alpha-rarefaction}. O valor de 45.000 normalizou o esforço amostral, garantindo que nenhuma das oito amostras fosse excluída das análises de diversidade e mitigou as disparidades na profundidade de sequenciamento. A diversidade alfa (riqueza de ASVs e índice de Shannon) foi calculada pelo comando \cmd{qiime diversity core-metrics}. A diferença entre as medianas dos dois ambientes foi avaliada pelo teste não-paramétrico de Kruskal-Wallis. A diferença na composição de comunidades (diversidade beta) foi estimada através das matrizes de distância de Jaccard, ordenadas por Análise de Coordenadas Principais (PCoA). A diferença entre os ambientes foi testada por Análise de Variância Molecular Permutacional (PERMANOVA) com 999 permutações.
